\chapter{Resumen}

El alumnado con discapacidad auditiva depende de que alguien transcriba lo que se dice en
clase. Este trabajo evalúa si las técnicas habituales de adaptación de modelos de
reconocimiento automático del habla mejoran esa transcripción en español, y desarrolla una
aplicación de subtitulado en vivo desplegable en la red de un centro. Se comparan tres
técnicas sobre Whisper con diseño pareado, intervalos de confianza por remuestreo y test de
signos: prompting contextual, post-corrección con un modelo de lenguaje local y ajuste fino
con LoRA. Solo la última mejora de forma concluyente, con 1,23 puntos de reducción del
error de palabra; la post-corrección lo empeora. El resultado principal es otro: segmentar
el audio por silencios en lugar de por reloj reduce el error 7,4 puntos a igual latencia,
seis veces más que la mejor técnica de adaptación. Se muestra además que el WER agregado
oculta los fallos que importan en accesibilidad, y se propone una métrica complementaria
que los expone.

\vspace{0.5cm}
\noindent\textbf{Palabras clave:} accesibilidad educativa; reconocimiento automático del
habla; adaptación al dominio; subtitulado en tiempo real; evaluación de sistemas ASR.

\chapter{Abstract}

Students with hearing impairment depend on someone transcribing what is said in class.
This work evaluates whether the usual domain adaptation techniques for automatic speech
recognition improve that transcription in Spanish, and develops a live captioning
application deployable on a school network. Three techniques are compared on Whisper using
a paired design, bootstrap confidence intervals and sign tests: contextual prompting,
post-correction with a local language model, and LoRA fine tuning. Only the last one
improves conclusively, reducing word error rate by 1.23 points; post-correction makes it
worse. The main finding lies elsewhere: segmenting audio at silences rather than at fixed
intervals reduces error by 7.4 points at equal latency, six times more than the best
adaptation technique. The work also shows that aggregate WER hides the failures that
matter for accessibility, and proposes a complementary metric that exposes them.

\vspace{0.5cm}
\noindent\textbf{Keywords:} educational accessibility; automatic speech recognition;
domain adaptation; real time captioning; ASR evaluation.
